% !TeX root = ../main.tex
% 第 1 章 拓扑群 （原书 CHAPTER 1, Topological groups, 印刷页 5-11）

\chapter{拓扑群}
\label{ch:1}

本章引入拓扑群的概念及其基本性质. 在无限扩张的情形下推广 Galois 理论基本定理时，这些概念是重要的. 读者可在 \cite{pon66} 与 \cite{hus66} 中进一步了解相关内容.

\section{定义与例子}

\begin{definition}\label{c1-def-1-1}
若元素集合 $G$ 满足
\begin{enumerate}[label=(\arabic*),leftmargin=*]
    \item $G$ 是一个群；
    \item $G$ 是一个拓扑空间；
    \item 映射 $G\times G\to G:(a,b)\mapsto ab$ 连续，其中 $G\times G$ 赋予乘积拓扑；
    \item 映射 $g\to G:x\mapsto x^{-1}$ 连续\footnote{原文即如此（作 $g\to G$）；按上下文应理解为 $G\to G$.}.
\end{enumerate}
则称 $G$ 为 \newterm{topological-group}{拓扑群}{topological group}.
\end{definition}

条件 (3) 与 (4) 分别展开叙述如下：
\begin{itemize}
    \item 若 $a,b\in G$，则对元素 $ab$ 的任一邻域 $W$，存在元素 $a$ 与 $b$ 的邻域 $U$ 与 $V$，使得 $UV\subset W$.
    \item 若 $a\in G$，则对元素 $a^{-1}$ 的任一邻域 $V$，存在元素 $a$ 的邻域 $U$，使得 $U^{-1}\subset V$.
\end{itemize}
最后，这两个条件可以合并表述为：
\begin{itemize}
    \item 若 $a,b\in G$，则对元素 $ab^{-1}$ 的任一邻域 $W$，存在元素 $a$ 与 $b$ 的邻域 $U$ 与 $V$，使得 $UV^{-1}\subset W$.
\end{itemize}

\begin{notation}\label{c1-not-1-2}
在本项目中，$G$ 的单位元一律记作 $1$.
\end{notation}

下面给出一些拓扑群的例子.

\begin{examples}\label{c1-ex-1-3}
以下是一些拓扑群的例子.
\begin{itemize}
    \item 若群 $G$ 赋予 \newterm{discrete-topology}{离散拓扑}{discrete topology}(即 $G$ 的每个子集都是开集)，则 $G$ 是拓扑群.
    \item 考虑带加法的实数群 $(\mathbb{R},+)$. 在 $\mathbb{R}$ 上赋予通常拓扑，它构成拓扑群.
    \item 更一般地，任何拓扑向量空间（例如 Banach 空间或 Hilbert 空间）的加法群都是拓扑群.
    \item 全体可逆 $n\times n$ 实矩阵构成的线性群 $GL_{n}(\mathbb{R})$，若赋予把 $GL_{n}(\mathbb{R})$ 视为 $\mathbb{R}^{n\times n}$ 的子集所得到的拓扑，则它是拓扑群.
\end{itemize}
\end{examples}

\begin{definition}\label{c1-def-1-4}
拓扑群的 \newterm{homomorphism-topgrp}{同态}{homomorphism} 是指连续的群同态. 拓扑群的 \newterm{isomorphism-topgrp}{同构}{isomorphism} 是指既是群同构又是同胚的映射.
\end{definition}

\section{初等性质}

拓扑群是均匀的，这意味着后文某些性质只需在单位元附近验证即可.

\begin{proposition}\label{c1-prop-2-1}
设 $G$ 是拓扑群，$a\in G$，映射 $f,g,\varphi:G\to G$ 定义为 $f(x)=xa$，$g(x)=ax$，$\varphi(x)=x^{-1}$，则 $f$、$g$ 与 $\varphi$ 都是 $G$ 的 \newterm{homeomorphism}{同胚}{homeomorphism}.
\end{proposition}

\begin{proof}
这些映射显然是双射；又由拓扑群的公理，它们是连续的.
\end{proof}

\begin{proposition}\label{c1-prop-2-2}
每个拓扑群都是均匀的，即对任意 $p,q\in G$，存在同胚 $f$ 使得 $f(p)=q$.
\end{proposition}

\begin{proof}
设 $p,q\in G$，令 $f:x\mapsto qp^{-1}x$. 由\autoref{c1-prop-2-1}，它是同胚.
\end{proof}

\begin{proposition}\label{c1-prop-2-3}
设 $V$ 是 $a$ 的一个邻域，则存在 $1$ 的邻域 $U$ 使得 $V=aU$.
\end{proposition}

\begin{proof}
定义 $U:=a^{-1}V$. 作为 $V$ 在连续映射 $g:G\to G:x\mapsto ax$ （见\autoref{c1-prop-2-1}）下的原像，$U$ 是 $1$ 的邻域. 而我们有 $V=aU$，结论得证.
\end{proof}

\begin{definition}\label{c1-def-2-1}
设 $G$ 是拓扑群，称邻域 $V$ 是 \newterm{symmetric-neighborhood}{对称的}{symmetric}，若 $V=V^{-1}$.
\end{definition}

\begin{proposition}\label{c1-prop-2-4}
设 $G$ 是拓扑群. 则 $G$ 存在由对称开集构成的拓扑基. 特别地，$G$ 的每个点处都有由对称邻域构成的\newterm{fundamental-system}{邻域基本系}{fundamental system of neighborhoods}.
\end{proposition}

\begin{proof}
设 $U$ 是 $G$ 的开子集，则 $U^{-1}$ 也是开集(它是 $U$ 在某连续映射下的原像)，于是 $V:=U\cap U^{-1}$ 也是开集. 显然有 $V\subseteq U$ 且 $V=V^{-1}$.
\end{proof}

\begin{proposition}\label{c1-prop-2-5}
设 $G$ 是拓扑群，$U$ 是 $x\in G$ 的一个邻域. 则：
\begin{enumerate}[label=(\arabic*),leftmargin=*]
    \item 存在 $x$ 的邻域 $V$，使得 $V^{2}\subseteq U$.
    \item 存在 $x$ 的邻域 $V$，使得 $V^{-1}\subseteq U$.
    \item 存在 $x$ 的邻域 $V$，使得 $VV^{-1}\subseteq U$.
    \item 存在 $x$ 的邻域 $V$，使得 $V^{-1}V\subseteq U$.
\end{enumerate}
\end{proposition}

\begin{proof}
由拓扑群的均匀性，只需对 $1$ 的邻域证明本命题.

上述四条断言都是下列运算连续性的推论：
\begin{enumerate}[label=(\arabic*),leftmargin=*]
    \item 考虑映射 $f_{1}:G\to G:x\mapsto x^{2}$，它连续（因为乘法连续），并令 $V=f_{1}^{-1}(U)$. 由 $1=1^{2}$，$V$ 是 $1$ 的邻域；并且 $V^{2}=f_{1}(V)=f_{1}(f_{1}^{-1}(U))\subseteq U$.
    \item 同理，利用 $f_{2}:G\to G:x\mapsto x^{-1}$ 与 $1=1^{-1}$.
    \item 同理，利用 $f_{3}:G\to G:x\mapsto xx^{-1}$ 与 $1=11^{-1}$.
    \item 同理，利用 $f_{4}:G\to G:x\mapsto x^{-1}x$ 与 $1=1^{-1}1$.
\end{enumerate}
\end{proof}

\begin{notation}\label{c1-not-2-2}
设 $X$ 是拓扑空间，$A\subseteq X$. 以 $\overline{A}$ 记 $A$ 的拓扑闭包.
\end{notation}

\begin{lemma}\label{c1-lem-2-6}
设 $X,Y$ 是两个拓扑空间，$f:X\to Y$ 是连续映射. 则 $f(\overline{A})\subseteq \overline{f(A)}$.
\end{lemma}

\begin{proof}
设 $y\in f(\overline{A})$，$W$ 是 $Y$ 中 $y$ 的一个邻域. 存在 $x\in\overline{A}$ 使得 $f(x)=y$. 由于 $f$ 连续，$U:=f^{-1}(W)$ 是 $x$ 的邻域. 于是 $U\cap A\neq\emptyset$ (因为 $x\in\overline{A}$). 从而
\[
\emptyset\neq f(U\cap A)\subseteq f(U)\cap f(A)\subseteq W\cap \overline{f(A)}.
\]
所以 $y\in\overline{f(A)}$.
\end{proof}

\begin{proposition}\label{c1-prop-2-7}
设 $G$ 是拓扑群，$\mathcal{V}$ 是 $1$ 的开邻域的一个基本系，$D$ 是 $G$ 中的稠密子集. 则 $\mathcal{B}:=\{Ux\mid x\in D,\,U\in\mathcal{V}\}$ 是 $G$ 的拓扑的一个基.
\end{proposition}

\begin{proof}
设 $W$ 是 $G$ 的开子集，$a\in W$. 则 $Wa^{-1}$ 是开集(它是 $W$ 在某连续映射下的原像)，且 $1\in Wa^{-1}$. 于是存在 $U\in\mathcal{V}$ 使得 $UU^{-1}\subseteq Wa^{-1}$ (利用 $\mathcal{V}$ 是基本系以及\autoref{c1-prop-2-5}).

由于 $D$ 稠密，$aD^{-1}$ 也稠密. 事实上，考虑连续（且双射的）映射 $f:G\to G:x\mapsto ax^{-1}$，有 $X=f(X)=f(\overline{D})\subseteq\overline{f(D)}$\footnote{原文即如此作 $X$；按上下文应为 $G$.} （见\autoref{c1-lem-2-6}）. 于是存在 $d\in U\cap aD^{-1}$，从而存在 $m\in D$ 使得 $d=am^{-1}$，这就推出 $m=d^{-1}a$. 从而 $Ud^{-1}a\in\mathcal{B}$.

由于 $d\in U$，我们得到 $Ud^{-1}a\subseteq UU^{-1}a\subseteq Wa^{-1}a=W$；这对一切 $a\in W$ 成立. 于是
\[
\bigcup_{a\in W}Ud^{-1}a\subseteq W.
\]
此外，$1\in Ud^{-1}$ (因为 $d\in U$)，故对一切 $a\in W$ 有 $a\in Ud^{-1}a$. 于是
\[
W\subseteq\bigcup_{a\in W}Ud^{-1}a.
\]
所以 $W=\bigcup_{a\in W}Ud^{-1}a$.
\end{proof}

\begin{proposition}\label{c1-prop-2-8}
设 $G$ 是拓扑群，$\mathcal{V}$ 是 $1$ 的开邻域的一个基本系. 则 $\bigcap_{V\in\mathcal{V}}V=\overline{\{1\}}$.
\end{proposition}

\begin{proof}
证明双向包含：

$\subseteq$：设 $x\in\bigcap_{V\in\mathcal{V}}V$，$W$ 是 $x$ 的一个邻域. 结合\autoref{c1-prop-2-3} 与\autoref{c1-prop-2-4}，知存在对称集 $V\in\mathcal{V}$ 使得 $Vx\subseteq W$. 由 $x$ 的假设，$x\in V=V^{-1}$，从而 $x^{-1}\in V$. 于是 $1=x^{-1}x\in Vx$. 故 $1\in U$，即 $U\cap\overline{\{1\}}\neq\emptyset$\footnote{原文此处即如此作 $U$；按上下文，此邻域应为上文所设的 $W$.}. 从而 $x\in\overline{\{1\}}$.

$\supseteq$：设 $x\in\overline{\{1\}}$，$V\in\mathcal{V}$ 对称(由\autoref{c1-prop-2-4}，可设 $V$ 对称). 由于 $xV$ 是 $x$ 的邻域，$xV\cap 1\neq\emptyset$. 这说明存在 $v\in V$ 使得 $xv=1$，于是 $x=v^{-1}\in V^{-1}=V$. 故 $x\in\bigcap_{V\in\mathcal{V}}V$.
\end{proof}

\begin{proposition}\label{c1-prop-2-9}
设 $G$ 是拓扑群，$A,B\subseteq G$. 则
\begin{enumerate}[label=(\arabic*),leftmargin=*]
    \item $\overline{A}\,\overline{B}=\overline{AB}$.
    \item $\overline{A}^{-1}=\overline{A^{-1}}$.
    \item $x\overline{A}y=\overline{xAy}$ 对一切 $x,y\in G$ 成立.
\end{enumerate}
\end{proposition}

\begin{proof}
(1) 设 $x\in\overline{A}$，$x\in\overline{B}$\footnote{原文即如此（作 $x\in\overline{B}$）；按上下文应为 $y\in\overline{B}$.}，$U$ 是 $1$ 的邻域. 存在 $1$ 的邻域 $V$ 使得 $xVyV\subseteq xyU$ (因为对一切 $b\in G$，映射 $G\to G:a\mapsto ba$ 连续).

由于 $xV$、$yV$ 分别是 $x$、$y$ 的邻域，存在 $a\in A\cap xV$ 与 $b\in B\cap yV$. 于是 $ab\in AB\cap xyU$，即 $xy\in\overline{AB}$.

(2) 证明双向包含.

$\subseteq$：对连续映射 $i:G\to G:z\mapsto z^{-1}$ 用\autoref{c1-lem-2-6}.

$\supseteq$：$A^{-1}\subseteq\overline{A^{-1}}\Rightarrow A\subseteq(\overline{A^{-1}})^{-1}=i(\overline{A^{-1}})$，而后一个集合是闭集，因为 $i$ 是同胚.

于是 $\overline{A}\subseteq(\overline{A^{-1}})^{-1}\Rightarrow(\overline{A})^{-1}\subseteq\overline{A^{-1}}$.

(3) 对同胚 $G\to G:z\mapsto xzy$ 同理.
\end{proof}

\section{拓扑的刻画}

利用拓扑群的均匀性，我们可以用单位元的邻域基来刻画拓扑.

\begin{proposition}\label{c1-prop-3-1}
设 $G$ 是拓扑群，$\mathcal{V}$ 是 $G$ 的单位元 $1$ 的一个邻域基. 则
\begin{enumerate}[label=(\arabic*),leftmargin=*]
    \item 对一切 $V_{1},V_{2}\in\mathcal{V}$，存在 $V\in\mathcal{V}$ 使得 $1\in V\subseteq V_{1}\cap V_{2}$.
    \item 对一切 $V\in\mathcal{V}$，存在 $W\in\mathcal{V}$ 使得 $W^{2}\subseteq V$.
    \item 对一切 $V\in\mathcal{V}$，存在 $W\in\mathcal{V}$ 使得 $W^{-1}\subseteq V$.
    \item 对一切 $V\in\mathcal{V}$ 与一切 $g\in G$，存在 $W\in\mathcal{V}$ 使得 $W\subseteq gVg^{-1}$.
    \item 对一切 $g\in G$，$\{gV\mid V\in\mathcal{V}\}$ 是 $g$ 的一个邻域基.
\end{enumerate}
反之，若 $G$ 是群，$\mathcal{V}$ 是 $G$ 的一个非空子集族且满足 (1)--(4)，则在 $G$ 上存在（唯一的）拓扑使 (5) 成立.
\end{proposition}

\begin{proof}
设 $\mathcal{V}$ 是拓扑群中 $1$ 的邻域基. 则 (1) 是邻域基定义的推论；(2)、(3)、(4) 是\autoref{c1-prop-2-5} 的推论；而 (5) 成立是因为映射 $f:G\to G:x\mapsto gx$ 是同胚.

反之，设 $G$ 是群，$\mathcal{V}$ 是 $G$ 的一个非空子集族且满足 (1)--(4). 注意 (1) 蕴含每个 $V\in\mathcal{V}$ 都含 $1$. 定义
\[
\mathcal{T}=\{U\subseteq G\mid \forall g\in U,\ \exists V\in\mathcal{V}\ \text{使得}\ gV\subseteq U\},
\]
并证明 $\mathcal{T}$ 是 $G$ 上的拓扑. 首先，空集与 $G$ 显然属于 $\mathcal{T}$. 设 $(U_i)_{i\in I}$ 是 $\mathcal{T}$ 的一族元素($I$ 为指标集)，并设 $g\in\bigcup_{i\in I}U_i$. 于是存在 $j\in I$ 使得 $g\in U_j$，由定义存在 $V\in\mathcal{V}$ 使得 $gV\subseteq U_j$，从而 $gV\subseteq\bigcup_{i\in I}U_i$. 再设 $U_{1},U_{2}\in\mathcal{T}$，$g\in U_{1}\cap U_{2}$. 由定义存在 $V_{1},V_{2}\in\mathcal{V}$ 使得 $gV_{1}\subseteq U_{1}$ 且 $gV_{2}\subseteq U_{2}$. 应用 (1) 即得 $V\in\mathcal{V}$ 使得 $V\subseteq V_{1}\cap V_{2}$，这意味着 $gV\subseteq U_{1}\cap U_{2}$. 故 $U_{1}\cap U_{2}\in\mathcal{T}$. 这样我们证明了 $\mathcal{T}$ 给出 $G$ 上的一个拓扑. 容易验证 (5) 成立.

下面证明 $G\to G:g\mapsto g^{-1}$ 连续. 设 $g\in G$，$U\in\mathcal{T}$ （即为开集）且 $g^{-1}\in U$. 我们要找到 $V\in\mathcal{V}$ 使得 $gV\subseteq U^{-1}$. 由 $\mathcal{T}$ 的定义，存在 $V\in\mathcal{V}$ 使得 $g^{-1}V\subseteq U$，从而 $V^{-1}g\subseteq U^{-1}$. 利用 (3) 可得 $V'\in\mathcal{V}$ 使得 $gV'\subseteq U^{-1}$；再利用 (4)，即知存在 $V''\in\mathcal{V}$ 使得 $gV''\subseteq g(g^{-1}Vg)\subseteq U^{-1}$.

再验证乘法 $G\times G\to G:(g_{1},g_{2})\mapsto g_{1}g_{2}$ 连续. 注意对一切 $V_{1},V_{2}\in\mathcal{V}$，集合 $g_{1}V_{1}\times g_{2}V_{2}$ 构成 $G\times G$ 中 $(g_{1},g_{2})$ 的邻域基. 设 $(g_{1},g_{2})\in G\times G$，$U$ 是满足 $g_{1}g_{2}\in U$ 的开集. 我们要找 $V_{1},V_{2}\in\mathcal{V}$ 使得 $g_{1}V_{1}g_{2}V_{2}\subseteq U$. 由于 $U\in\mathcal{T}$，存在 $V\in\mathcal{V}$ 使得 $g_{1}g_{2}V\subseteq U$. 用 (2) 得 $V'\in\mathcal{V}$ 使得 $g_{1}g_{2}V'V'\subseteq U$. 再用 (4)，得 $V''\in\mathcal{V}$ 使得 $V''\subseteq g_{2}V'g_{2}^{-1}$，从而 $g_{1}V''g_{2}V'\subseteq g_{1}(g_{2}V'g_{2}^{-1})g_{2}V'\subseteq U$. 证明完毕.
\end{proof}

\section{拓扑子群}

本节简要讨论拓扑群的一些性质，特别是其子群的性质.

\begin{proposition}\label{c1-prop-4-1}
设 $G$ 是拓扑群，$H$ 是子群. 则 $H$ 是拓扑群.
\end{proposition}

\begin{proof}
只需证明 $H$ 上的乘法与取逆连续，而这显然是对的：它们只是 $G$ 上相应映射到 $H$ 上的限制.
\end{proof}

\begin{proposition}\label{c1-prop-4-2}
设 $G$ 是拓扑群，$H$ 是子群. 则 $\overline{H}$ 是 $G$ 的子群. 此外，若 $H$ 正规，则 $\overline{H}$ 也正规.
\end{proposition}

\begin{proof}
只需证明 $\overline{H}$ 对乘法与取逆封闭\footnote{原文此处用词为 addition（加法）；按上下文应为乘法.}. 设 $a,b\in\overline{H}$，证明 $ab^{-1}\in\overline{H}$.

设 $W$ 是 $ab^{-1}$ 的邻域，则存在 $a$ 与 $b$ 的邻域 $U$ 与 $V$ 使得 $UV^{-1}\subseteq W$ (因为映射 $G\times G\to G:(x,y)\mapsto xy^{-1}$ 连续). 我们有 $U\cap H\neq\emptyset$ 且 $V\cap H\neq\emptyset$ (注意 $a,b\in\overline{H}$). 于是可取 $x\in U\cap H$ 与 $y\in V\cap H$，使得 $xy^{-1}\in UV^{-1}\cap H\subseteq W\cap H$. 故 $W\cap H\neq\emptyset$，从而 $ab^{-1}\in\overline{H}$.

再设 $H$ 正规. 取 $a\in\overline{H}$ 与 $g\in G$，我们证明 $g^{-1}ag\in\overline{H}$，这便蕴含结论.

设 $V$ 是 $g^{-1}ag$ 的邻域，则存在 $a$ 的邻域 $U$ 使得 $g^{-1}Ug\subseteq V$. 由 $a\in\overline{H}$ 知 $U\cap H\neq\emptyset$. 于是存在 $h\in U\cap H$，使得 $g^{-1}hg\in V$. 又因 $H$ 正规，$g^{-1}hg\in H$. 故 $V\cap H\neq\emptyset$.
\end{proof}

\begin{definition}\label{c1-def-4-1}
称拓扑空间是 \newterm{totally-disconnected}{全不连通的}{totally disconnected}，若它的连通子空间只有单点集.
\end{definition}

\begin{examples}\label{c1-ex-4-2}
\
\begin{itemize}
    \item 一切离散空间都是全不连通的.
    \item 有理数集 $\mathbb{Q}$ 与无理数集 $\mathbb{R}\setminus\mathbb{Q}$ 都是全不连通空间.
    \item 实数集 $\mathbb{R}$（赋予通常拓扑）不是全不连通空间.
\end{itemize}
\end{examples}

\begin{lemma}\label{c1-lem-4-3}
设 $X$ 是全不连通拓扑空间. 则对每个 $x\in X$，$\{x\}$ 在 $X$ 中是闭集.
\end{lemma}

\begin{proof}
设 $C$ 是 $\{x\}$ 的拓扑闭包. 若 $A|B$ 是 $C$ 的一个分离且 $x\in A$，则 $A$ 在 $C$ 中闭，从而在 $X$ 中闭. 于是 $A=C$. 这样 $C$ 是连通的，而 $X$ 全不连通，故必有 $C=\{x\}$.
\end{proof}

\begin{proposition}\label{c1-prop-4-4}
设 $G$ 是拓扑群.
\begin{enumerate}[label=(\arabic*),leftmargin=*]
    \item 若 $H$ 是 $G$ 的开（相应地：闭）子群，则每个陪集 $Hg$ 或 $gH$ 都是开（相应地：闭）集.
    \item $G$ 的每个开子群都是闭的. 若 $G$ 紧致，则 $G$ 的每个开子群的指数有限.
    \item $G$ 的每个指数有限的闭子群都是开的.
    \item 若 $H$ 是包含 $G$ 的某个非空开子集 $U$ 的子群，则 $H$ 在 $G$ 中是开集.
    \item $G$ 是 Hausdorff 空间当且仅当 $\{1\}$ 是 $G$ 的闭集. 且若 $K$ 是 $G$ 的正规子群，则 $G/K$ 是 Hausdorff 空间当且仅当 $K$ 在 $G$ 中闭. 若 $G$ 全不连通，则 $G$ 是 Hausdorff 空间.
\end{enumerate}
\end{proposition}

\begin{proof}
(1) 这是\autoref{c1-prop-2-2} 的推论.

(2) 设 $H$ 开，证明 $\overline{H}\subseteq H$. 取 $a\in\overline{H}$. 由于 $H$ 开，$aH$ 是 $a$ 的开邻域，故 $aH\cap H\neq\emptyset$. 于是存在 $h_{1},h_{2}\in H$ 使得 $ah_{1}=h_{2}$，即 $a=h_{2}h_{1}^{-1}\in H$. 至于 $H$ 指数有限：注意各 $gH$ 是开集（由 (1)）、两两不交且并为 $G$. 故若 $G$ 紧致，$H$ 必指数有限.

(3) 设 $H$ 是 $G$ 的指数有限的闭子群. 则对每个 $g\in G$，$gH$ 闭. 而 $G\setminus H=\bigcup_{g\notin H}gH$，$H$ 指数有限蕴含 $G\setminus H$ 是闭集（闭集的有限并），故 $H$ 是开集.

(4) 每个集合 $hU$ 都是开集（由 (1)），且 $H=\bigcup_{h\in H}hU$，故 $H$ 开.

(5) 由一般拓扑知识，Hausdorff 空间中的单点集都是闭集. 反之设 $\{1\}$ 在 $G$ 中闭. 取 $G$ 中不同元素 $a,b$. 由\autoref{c1-prop-2-2}，集合 $\{ab^{-1}\}$ 是闭集(它是 $\{1\}$ 在同胚下的像). 于是存在开集 $U$ 使得 $1\in U$ 且 $U\subseteq G\setminus\{ab^{-1}\}$. 由映射 $G\to G:(x,y)\mapsto xy^{-1}$ （显然）连续\footnote{原文即如此作 $G\to G$；按上下文应为 $G\times G\to G$.}，$U$ 的原像是开集. 于是存在含 $1$ 的开集 $V,W\subseteq G$，使得 $VW^{-1}\subseteq U$. 这样 $a^{-1}b\notin VW^{-1}$，从而 $aV\cap bW=\emptyset$. 由于 $aV$ 与 $bW$ 都是开集，结论得证.

后面的断言由第一条断言、商拓扑的定义与\autoref{c1-lem-4-3} 推出.
\end{proof}

\section{商群}

这里给出关于拓扑群的商的一个命题.

\begin{proposition}\label{c1-prop-5-1}
设 $G$ 是拓扑群，$N$ 是正规子群. 则 $G/N$ 赋予商拓扑后是拓扑群. 此外，典范投影 $\pi:G\to G/N$ 是连续的开同态.
\end{proposition}

\begin{remark}\label{c1-rem-5-1}
商拓扑定义为\footnote{原文即如此作 $U\subseteq X$；按上下文应为 $U\subseteq G$.}
\[
\mathcal{T}=\{U\subseteq X\mid \pi^{-1}(U)\ \text{在}\ G\ \text{中开}\}.
\]
\end{remark}

\begin{proof}
先证 $\pi$ 是开映射. 设 $U$ 是 $G$ 的开子集，则 $\pi^{-1}(\pi(U))=UN$，而 $UN$ 是开集(因为 $U$ 开)，故结论成立. 此外，由 $\mathcal{T}$ 的定义，$\pi$ 的连续性是平凡的.

下面证明 $G/N$ 是拓扑群. $N$ 正规蕴含 $G/N$ 是抽象群. 验证乘法与取逆连续.

设 $W$ 是 $g_{1}g_{2}N\in G/N$ 的邻域，则因 $\pi$ 连续，$\pi^{-1}(W)$ 是 $g_{1}g_{2}$ 的邻域. 由 $G$ 中乘法的连续性，存在 $g_{1},g_{2}$ 的邻域 $U_{1},U_{2}$，使得 $U_{1}U_{2}\subseteq\pi^{-1}(W)$. 于是 $\pi(U_{1})$ 与 $\pi(U_{2})$ 分别是 $g_{1}N$ 与 $g_{2}N$ 的邻域(利用 $\pi$ 是开映射). 又因 $\pi$ 是同态，$\pi(U_{1})\pi(U_{2})=\pi(U_{1}U_{2})\subseteq W$. 故乘法连续.

设 $W$ 是 $gN\in G/N$ 的邻域，则因 $\pi$ 连续，$\pi^{-1}(W)$ 是 $g$ 的邻域. 由 $G$ 中取逆的连续性，存在 $g$ 的邻域 $U$ 使得 $U^{-1}\subseteq\pi^{-1}(W)$. 于是 $\pi(U)$ 是 $gN$ 的邻域，且 $\pi(U)^{-1}=\pi(U^{-1})\subseteq W$. 故取逆连续.
\end{proof}

\section{Lie 群}

我们简要介绍 Lie 群，因为它们提供大量拓扑群的例子. 关于这一主题的更多信息可参见 ??\footnote{原文此处即为 ``??''（未填妥的引文）.}.

\begin{definition}\label{c1-def-6-1}
\newterm{Lie-group}{Lie 群}{Lie group} 是指赋予群结构的有限维光滑流形 $G$，其乘法光滑. 这就是说，存在光滑乘法 $\mu_{G}\times G\to G$\footnote{原文即如此排印；按上下文应理解为 $\mu:G\times G\to G$.}、取逆 $i:G\to G$ 与单位元 $1\in G$，满足群公理.
\end{definition}

光滑流形的定义可见 \cite[Chapter 2]{lang02b}.

\begin{remark}\label{c1-rem-6-2}
注意 Lie 群都是拓扑群. 它们为拓扑群提供了大量例子.
\end{remark}

\begin{examples}\label{c1-ex-6-3}
\
\begin{itemize}
    \item $\mathbb{R}$ 与 $\mathbb{C}$ 对加法构成 Lie 群. 并且，任何有限维实或复向量空间对加法都构成 Lie 群.
    \item $\mathbb{R}^{*}(=\mathbb{R}\setminus\{0\})$ 与 $\mathbb{C}^{*}$ 对乘法构成 Lie 群. $\mathbb{S}=\{z\in\mathbb{C}\mid |z|=1\}$ 对乘法也构成 Lie 群.
    \item 若 $G$ 与 $H$ 是 Lie 群，则乘积 $G\times H$ 是 Lie 群（赋予通常的乘积结构）. 于是 (1) 与 (2) 蕴含环面 $\mathbb{T}_{n}=\mathbb{S}^{n}$ 是 Lie 群.
    \item Lie 群的基本例子是有限维实向量空间 $V$ 上全体可逆线性映射构成的群 $GL(V)$. 我们对 $V=\mathbb{R}^{n}$ 的情形加以验证. 设 $f\in L(V,V)$，取标准基 $\{e_{1},\ldots,e_{n}\}$. 元素 $f(e_{i})\in\mathbb{R}^{n}$ 就是 $f$ 的矩阵的第 $i$ 列. 这给出 $L(V,V)$ 到 $M_{n}(\mathbb{R})=\mathbb{R}^{n^{2}}$ 的映射. 行列式定义光滑函数 $\det:M_{n}(\mathbb{R})\to\mathbb{R}$. 特别地，$GL_{n}(\mathbb{R})=\det^{-1}(\mathbb{R}_{n}^{*})$ 是 $\mathbb{R}^{2}_{n}$\footnote{原文即如此作 $\mathbb{R}^{2}_{n}$；按上下文应为 $\mathbb{R}^{n^{2}}$.} 的开子集，从而是光滑流形. 此外，两矩阵 $A$ 与 $B$ 之积的元素是 $A$ 与 $B$ 的元素的多项式，故乘法定义光滑映射.
    \item 特殊线性群 $SL_{n}(\mathbb{R})$ 定义为行列式映射 $\det:GL_{n}(\mathbb{R})\to\mathbb{R}^{*}$ 的核，它是 $GL_{n}(\mathbb{R})$ 的子群，故为 Lie 群.
    \item 经典群 $O_{n}(\mathbb{R})$ 与 $SO_{n}(\mathbb{R})$ 也是 Lie 群，其中 $O_{n}(\mathbb{R})=\{A\in GL_{n}(\mathbb{R})\mid AA^{T}=A^{T}A=\mathrm{Id}_{n}\}$，$SO_{n}(\mathbb{R})=\{A\in O_{n}(\mathbb{R})\mid \det(A)=1\}$.
\end{itemize}
\end{examples}

下面给出一些不是 Lie 群的拓扑群的例子.

\begin{examples}\label{c1-ex-6-4}
\
\begin{itemize}
    \item 无限维群，例如无限维实向量空间的加法群. 它们不是有限维流形，故不是 Lie 群.
    \item 某些全不连通群，例如域的无限扩张的 Galois 群，或 $p$ 进数域的加法群. 它们的底空间不是实流形，故不是 Lie 群.
\end{itemize}
\end{examples}
